\documentclass[10pt,letterpaper]{article}
\usepackage[margin=0.65in]{geometry}
\usepackage{booktabs}
\usepackage{longtable}
\usepackage{graphicx}
\usepackage{hyperref}
\title{Global Phylogenomic Radiation of Bemisia Whiteflies and Allied Aleyrodidae (Hemiptera: Sternorrhyncha), Native Rust Quartet Coalescent Species-Tree Inference, Symbiotic and Cytochrome P450 Selection, and Cross-Regional vs. Multi-Locus Demographic Skylines}
\author{Bryan P. White \\ CRV Science Labs}
\date{November 9, 2026 \quad|\quad Preprint ID: PREPRINT-23 (v1.0.0 Empirical Preprint) \quad|\quad DOI: \href{https://doi.org/10.5281/zenodo.24093136}{10.5281/zenodo.24093136}}
\begin{document}
\maketitle
\begin{abstract}
We reconstruct the Triassic-to-Anthropocene global phylogenomic radiation, symbiotic and metabolic enzyme evolution, and multi-compartment demographic history of Bemisia whiteflies and allied Aleyrodidae (Hemiptera: Sternorrhyncha) across 544 active verified NCBI sequence records (plus 6 autonomously quarantined GenBank UNVERIFIED anomalies) comprising a 64-taxon, 18-locus supermatrix (L = 90,311 aligned columns), 48 nuclear functional coding sequences, and 12 single-species regional and multi-locus Bemisia tabaci cohorts (12 x N = 36 = 432 isolates). Executed end-to-end via autonomous Native Rust bio-suite-supervisor optimization across 32 candidate RuNIE chains (total runtime 503.59 s; B = 1,000 UFBoot; 25,000 MCMC generations with 20\% burn-in), RY3 purine/pyrimidine recoding under GTR+F+R4 reduces Sternorrhynchan compositional heterogeneity by 88.0\% (RCFV = 0.0063 vs. 0.0525, ln L = -631,907.97), dating the Sternorrhyncha root to 234.6 Ma, Crown Aleyrodidae to 129.4 Ma, the Crown Bemisia tabaci cryptic species complex to 20.6 Ma, and the MED-MEAM1-IO crown divergence to 2.86 Ma. Native Rust quartet multispecies coalescent inference (runie-quartet-coalescent across 18 gene trees and 11,436,768 quartets in 89.77 ms) recovers 8,310,800 satisfied quartets (normalized Q = 0.7267, mean LPP = 0.8602) while ABBA-BABA tests quantify asymmetric MED-MEAM1 and SSA1-IO admixture. Sixty-one-state Goldman-Yang codon models identify strong purifying constraint on Orco chemosensation (omega\_0 = 0.208) and Portiera bacteriocyte symbiosis alongside diversifying selection on salivary BtFer1/ALP hydrolases (omega\_2 = 1.00) and CYP6CM1 substrate-recognition pockets (omega\_2 = 1.00). Finally, 12 coupled Bayesian coalescent skylines contrast Old World savanna and agrarian expansions against New World founder bottlenecks and reveal sharp locus-specific effective population compressions at nuclear CYP6CM1 and ChE coding loci relative to neutral mitochondrial COX1 (1Ne) and Katoh E-INS-i-aligned nuclear ITS1-5.8S (4Ne) reservoirs.
\end{abstract}
\section{Introduction}
We reconstruct the Triassic-to-Anthropocene global radiation of \textit{Bemisia} whiteflies and allied Aleyrodidae (Hemiptera: Sternorrhyncha) across 544 active verified NCBI sequence records (64 complete Sternorrhyncha/Aleyrodidae/\textit{Bemisia} mitogenomes, 48 nuclear functional CDS orthologs across 4 physiological gene families, and 12 single-species regional and multi-locus \textit{Bemisia tabaci} cohorts totaling 432 isolates; $L = 90,311$ aligned supermatrix columns) plus 6 autonomously quarantined GenBank UNVERIFIED anomalies using 100\% Native Rust tools in \texttt{bio-suite-supervisor} \cite{ref1,ref5,ref7}.
\section{Autonomous Pipeline Supervision and 32-Chain RuNIE Sweep}
Across 32 candidate Native Rust chains spanning four data-hygiene and recoding regimes, purine/pyrimidine \texttt{RY3} recoding under \texttt{GTR+F+R4} achieved $\ln L = -631,907.97$ and reduced compositional heterogeneity to $\mathrm{RCFV} = 0.0063$ in 503.59 s end-to-end execution.
\section{Calibrated Phylogenomics, Quartet Coalescent and Demographic Skylines}
Fossil-calibrated \texttt{TimeTreeLSD2} dating places the Sternorrhyncha root at 234.60 Ma, Crown Aleyrodidae at 129.40 Ma, Crown \textit{Bemisia tabaci} complex at 20.60 Ma, and Crown MED--MEAM1--IO at 2.86 Ma \cite{ref2,ref3,ref4}. Native 18-locus Quartet Multispecies Coalescent inference across 11,436,768 quartets recovers normalized $Q = 0.7267$ (mean LPP = 0.8602) \cite{ref17}, while 61-state Goldman--Yang codon models identify diversifying selection in \texttt{CYP6CM1} ($\omega_2 = 1.00$) \cite{ref16,ref19} and 12 coupled 25,000-generation Bayesian skylines quantify cross-regional vs. multi-locus coalescent discordance \cite{ref18}.
\begin{thebibliography}{99}
\bibitem{ref1} White, B. P. (2026). RuNIE (Rust Native Inference Engine): An Open-Source Architecture for Autonomous Bioinformatics Pipeline Synthesis, Blinded Validation, and Bayesian Scope-Depth Auditing. CRV Science Labs Preprint / Zenodo. DOI: 10.5281/zenodo.23969275 \url{https://doi.org/10.5281/zenodo.23969275}

\bibitem{ref2} Misof, B., Liu, S., Meusemann, K., Peters, R. S., Donath, A., Mayer, C., et al. (2014). Phylogenomics resolves the timing and pattern of insect evolution. Science, 346(6210), 763--767. DOI: 10.1126/science.1257570 \url{https://doi.org/10.1126/science.1257570}

\bibitem{ref3} Johnson, K. P., Dietrich, C. H., Friedrich, F., Beutel, R. G., Wipfler, B., Peters, R. S., et al. (2018). Phylogenomics and the evolution of hemipteroid insects. Proceedings of the National Academy of Sciences, 115(50), 12775--12780. DOI: 10.1073/pnas.1815820115 \url{https://doi.org/10.1073/pnas.1815820115}

\bibitem{ref4} Drohojowska, J., Szwedo, J.,  y a, D., Huang, D. Y., \& M ller, P. (2020). Fossils reshape the Sternorrhyncha evolutionary tree (Insecta, Hemiptera). Scientific Reports, 10, 11390. DOI: 10.1038/s41598-020-69029-4 \url{https://doi.org/10.1038/s41598-020-69029-4}

\bibitem{ref5} De Barro, P. J., Liu, S. S., Boykin, L. M., \& Dinsdale, A. B. (2011). Bemisia tabaci: A statement of species status. Annual Review of Entomology, 56, 1--19. DOI: 10.1146/annurev-ento-112408-085504 \url{https://doi.org/10.1146/annurev-ento-112408-085504}

\bibitem{ref6} Mugerwa, H., Seal, S., Wang, H. L., Patel, M. V., Kabaalu, R., Omongo, C. A., et al. (2018). African ancestry of New World, belonging to the Bemisia tabaci cryptic species complex, and morphological distinctions among Sub-Saharan lineages. PLoS ONE, 13(5), e0196733. DOI: 10.1371/journal.pone.0196733 \url{https://doi.org/10.1371/journal.pone.0196733}

\bibitem{ref7} Boykin, L. M., Bell, C. D., Evans, G., Small, I., \& De Barro, P. J. (2013). Is agriculture driving the diversification of the Bemisia tabaci species complex (Hemiptera: Sternorrhyncha: Aleyrodidae)? Dating, host-plant transitions and global biogeography. Molecular Phylogenetics and Evolution, 69(3), 1084--1092. DOI: 10.1016/j.ympev.2013.08.018 \url{https://doi.org/10.1016/j.ympev.2013.08.018}

\bibitem{ref8} Thao, M. L., Baumann, L., \& Baumann, P. (2004). Organization of the mitochondrial genomes of whiteflies, aphids, and psyllids (Hemiptera, Sternorrhyncha). BMC Evolutionary Biology, 4, 25. DOI: 10.1186/1471-2148-4-25 \url{https://doi.org/10.1186/1471-2148-4-25}

\bibitem{ref9} Tay, W. T., Elfekih, S., Polaszek, A., Evans, G. A., Gordon, K. H. J., \& De Barro, P. J. (2017). Novel molecular approach to define pest species status and taxonomic validity within the Bemisia tabaci cryptic species complex. Scientific Reports, 7, 489. DOI: 10.1038/s41598-017-00528-7 \url{https://doi.org/10.1038/s41598-017-00528-7}

\bibitem{ref10} Elfekih, S., Etter, P., Tay, W. T., Fumagalli, M., Gordon, K., Johnson, E., \& De Barro, P. J. (2018). Genome-wide analyses of the Bemisia tabaci species complex reveal contrasting patterns of admixture and evolutionary divergence across global lineages. BMC Genomics, 19, 252. DOI: 10.1186/s12864-018-4601-5 \url{https://doi.org/10.1186/s12864-018-4601-5}

\bibitem{ref11} Chen, W., Hasegawa, D. K., Kaur, N., Kliot, A., Pinheiro, P. V., Luan, J., et al. (2016). The draft genome of whitefly Bemisia tabaci MEAM1, a global crop-associated hemipteran, provides molecular insights into evolution and endosymbiosis. Proceedings of the National Academy of Sciences, 113(52), E8496--E8505. DOI: 10.1073/pnas.1610539113 \url{https://doi.org/10.1073/pnas.1610539113}

\bibitem{ref12} Xie, W., Chen, C., Yang, Z., Guo, L., Yang, X., Wang, D., et al. (2017). Genome sequencing of the Mediterranean species of the Bemisia tabaci complex provides insights into cytochrome P450 monooxygenase expansion and host adaptation. GigaScience, 6(5), gix018. DOI: 10.1093/gigascience/gix018 \url{https://doi.org/10.1093/gigascience/gix018}

\bibitem{ref13} Sloan, D. B., \& Moran, N. A. (2012). Endosymbiotic bacteria as a source of carotenoids and essential amino acids in whiteflies. Biology Letters, 8(6), 986--989. DOI: 10.1098/rsbl.2012.0664 \url{https://doi.org/10.1098/rsbl.2012.0664}

\bibitem{ref14} Luan, J. B., Chen, W., Hasegawa, D. K., Simmons, A. M., Wintermantel, W. M., Ling, K. S., et al. (2015). Metabolic coevolution in the bacterial symbiosis of whiteflies and related plant-sap-feeding insects. Genome Biology and Evolution, 7(9), 2635--2647. DOI: 10.1093/gbe/evv170 \url{https://doi.org/10.1093/gbe/evv170}

\bibitem{ref15} Su, Q., Peng, Z., Tong, H., Xie, W., Wang, S., Wu, Q., Zhang, J., Li, C., \& Zhang, Y. (2019). A salivary ferritin in the whitefly modulates host plant signaling and iron homeostasis during phloem feeding. Proceedings of the National Academy of Sciences, 116(24), 11695--11704. DOI: 10.1073/pnas.1905446116 \url{https://doi.org/10.1073/pnas.1905446116}

\bibitem{ref16} Karunker, I., Benting, J., Lueke, B., Ponge, T., Nauen, R., Roditakis, E., et al. (2008). Over-expression of cytochrome P450 CYP6CM1 is associated with metabolic monooxygenase activity across Mediterranean and Middle East--Asia Minor Bemisia tabaci lineages. Insect Biochemistry and Molecular Biology, 38(6), 634--644. DOI: 10.1016/j.ibmb.2008.03.008 \url{https://doi.org/10.1016/j.ibmb.2008.03.008}

\bibitem{ref17} Mirarab, S., \& Warnow, T. (2015). ASTRAL-II: Coalescent-based species tree estimation with many hundreds of taxa and thousands of genes. Bioinformatics, 31(12), i44--i52. DOI: 10.1093/bioinformatics/btv234 \url{https://doi.org/10.1093/bioinformatics/btv234}

\bibitem{ref18} Drummond, A. J., Rambaut, A., Shapiro, B., \& Pybus, O. G. (2005). Bayesian coalescent inference of past population dynamics from molecular sequences. Molecular Biology and Evolution, 22(5), 1185--1192. DOI: 10.1093/molbev/msi103 \url{https://doi.org/10.1093/molbev/msi103}

\bibitem{ref19} Yang, Z., Wong, W. S. W., \& Nielsen, R. (2005). Bayes empirical Bayes inference of amino acid sites under positive selection. Molecular Biology and Evolution, 22(4), 1107--1118. DOI: 10.1093/molbev/msi097 \url{https://doi.org/10.1093/molbev/msi097}

\end{thebibliography}

\begin{figure}[htbp]
\centering
\includegraphics[width=0.98\textwidth]{figures/fig6-runie-pipeline-architecture.pdf}
\caption{Dual Empirical RuNIE Pipeline Architectures, Full-Length Alignments \& 4-Rank Combinatorial Model Comparison Across Pipeline A (64-Taxon, 18-Locus Supermatrix) \& Pipeline B (12 Single-Species Regional \& Multi-Locus Coalescent Skyline Cohorts, $N = 432$ Isolates).}
\label{fig:fig6}
\end{figure}

\end{document}
